\documentclass{article}
\usepackage{fullpage}
\begin{document}
\noindent\textbf{Describe the type(s) of research subjects or activities that interest you at your first and second choice host laboratories, and discuss any particular factors influencing your choice of host laboratories.}
\\\\

I am interested in research that models complex, large-scale physical systems using molecular dynamics
and quantum algorithms. I’m seeking environments where I can deepen my understanding of large-scale
simulation, hybrid quantum-classical computing, and optimization methods to solve scientific problems in
material sciences, chemistry, and biology. Oak Ridge National Laboratory (ORNL), my first choice, and Brookhaven National Laboratory (BNL), my second choice, strongly align with these goals through their leadership in classical and quantum computing.
\\

At ORNL, I aim to build on my experience with hybrid quantum-classical workflows and the molecular
dynamics of biomolecular systems. My current work involves group-theoretic methods to minimize the
resources needed for the classical simulation of quantum circuits, and I’m interested in learning more about
other hybrid quantum-classical approaches for applications in chemistry, materials science, and combinatorial
optimization. I am particularly excited about working in the Quantum Computational Science Group under Dr. Ryan
Bennink and their Quantum High-Performance Computing (Q-HPC) initiative, which focuses on integrating
quantum algorithms with exascale classical computing. My background in quantum algorithms, high-performance computing, molecular simulation, and Monte Carlo optimization provides a strong foundation to
contribute to this initiative. My work as part of the MSU-DOE Plant Research Laboratory has often involved modelling large
biomolecular systems. These are expensive to simulate for longer timescales, which limits the range of
biological phenomena that can be studied. Through the Molecular Biophysics Group and Oak Ridge Leadership Computing Facility at ORNL, I hope to learn how large-scale molecular dynamics and deep-learning-based structure prediction methods can give new insights into biomolecular processes.
\\

At BNL, I’m drawn towards the Computing and Data Sciences Directorate and their focus on quantum
algorithms for quantum chromodynamics simulation. My previous experience has involved exploiting physical
structure to reduce the cost of quantum computation, and I am eager to learn how such
optimizations can be applied in high energy physics. I am also interested in learning about exascale numerical
algorithms for solving problems in materials science and biology as part of the Computational Science Initiative
at BNL. By doing so, I hope to contribute to HPC workflow engineering and quantum circuit design.


\end{document}